\documentclass[11pt]{article}
\usepackage[a4paper,margin=26mm]{geometry}
\usepackage[T1]{fontenc}
\usepackage{lmodern,amsmath,amssymb,amsthm,booktabs,array,longtable}
\usepackage[colorlinks=true,linkcolor=blue,citecolor=blue,urlcolor=blue]{hyperref}
\newtheorem{theorem}{Theorem}
\newtheorem{corollary}[theorem]{Corollary}
\newcommand{\R}{\mathbb R}
\newcommand{\sgn}{\operatorname{sgn}}
\newcommand{\Pf}{\operatorname{Pf}}
\title{A cubic counterexample at Schnell--Ketzmerick--Eckardt\
(2018, Sec.\ V, p.\ 11, Eq.\ (55)),\
and an extension of their Ohmic bound}
\author{Haobo Ma$^{1}$ \and Wenlin Zhang$^{2,1}$}
\date{5 November 2026\\[3pt]
\small $^1$ChronoAI Pte Ltd, Singapore\\
\small $^2$National University of Singapore}
\begin{document}
\maketitle
\begin{abstract}
For the non-Ohmic expectation $M_S\leq 2M_B-1$ of
Schnell, Ketzmerick and Eckardt (2018, arXiv:1802.04088v2,
Sec.\ V, p.\ 11, Eq.\ (55) and the surrounding paragraph),
we give a cubic counterexample and extend their Ohmic bound
to odd monomial spectra.
Two baths with $J(x)=x^3$ have five Bose-selected modes.
The normalized coupling weights are $(5,4,4,4,4)/21$ and
$(6,6,3,3,3)/21$ at energies $(1,2,3,4,5)$, and the selected
occupation fractions are $(14,56,21,8,1)/100$.
Positive signed Pfaffian minors show that this selection is unique
and persists on an open set of parameters, including spectral
exponents near three. Five selected modes persist after adding
any finite number of sufficiently weakly coupled higher levels.
For odd $q$, we prove $M_S\leq(q+1)M_B-1$ when $J(x)=x^q$
and the selected principal matrix has a one-dimensional kernel;
$q=1$ recovers the source's Ohmic bound, while $q=3$ gives seven
as an upper bound for two baths under this hypothesis.
\end{abstract}

\section{The expectation and the selection condition}
For two product-coupled reservoirs, Schnell, Ketzmerick and
Eckardt~\cite[Sec.\ V, Eq.\ (50)]{ske} use the rate-asymmetry matrix
\begin{equation}\label{eq:matrix}
 A^{(d)}_{ij}=J_d(\varepsilon_j-\varepsilon_i)(f_i f_j-g_i g_j),
 \qquad J_d(x)=|x|^d\sgn(x),\quad J_d(0)=0,
\end{equation}
where $d>0$, $f_i,g_i>0$, and $\sum_i f_i=\sum_i g_i=1$.
A common positive prefactor, such as $2\pi/\hbar$, does not
change selection and is suppressed. These weights correspond to
unit vectors $|f\rangle=\sum_i\sqrt{f_i}|i\rangle$ and
$|g\rangle=\sum_i\sqrt{g_i}|i\rangle$, so their rank-one
projectors have exactly the coupling form in the source.

A nonempty set $S\subseteq\{1,\ldots,M\}$ is selected if there
is a vector $\eta$ satisfying
\begin{equation}\label{eq:selection}
 (A\eta)_i=0,\ \eta_i>0\quad(i\in S),\qquad
 (A\eta)_i<0,\ \eta_i=0\quad(i\notin S).
\end{equation}
This is Eq.\ (11) of~\cite{ske}; we fix its arbitrary scale by
$\sum_i\eta_i=1$. Write $M_S=|S|$ and let $M_B$ count the baths.
For Ohmic spectra ($d=1$), the source gives a generic linear
independence argument for $M_S\leq2M_B-1$. Immediately around
Eq.\ (55), p.\ 11, it says ``We expect that there is a generalization''
to non-Ohmic spectra and, after that equation, says ``similarly
observe $M_S\leq2M_B-1$'' for $J_d$.
Our counterexample answers the proposed extension negatively;
it does not change what occurred in that finite numerical sample.
The expectation is repeated in Schnell's thesis,
Sec.\ 4.4, p.\ 80, Eq.\ (4.54)~\cite{schnell}.
Our example concerns this non-Ohmic extrapolation.

The signed-Pfaffian description of kernels of odd antisymmetric
matrices is standard in this setting; see
Knebel, Weber, Kr\"uger and Frey~\cite[Methods, Eq.\ (19)]{knebel}.
Geiger, Knebel and Frey~\cite[Sec.\ 2.2]{geiger} relate positive
kernels to coexistence in antisymmetric Lotka--Volterra systems.
Here we give the five-dimensional identity explicitly, so the
counterexample requires only finite arithmetic and continuity.
Later work by Petiziol and Eckardt~\cite[Eq.\ (6), Fig.\ 3]{petiziol}
engineers five selected modes using frequency-selective artificial
reservoirs. Its Lorentzian spectra differ from the two cubic
product-coupled reservoirs considered here.

\section{An exact five-mode example}
\begin{theorem}\label{thm:example}
For $M=5$, $d=3$, and
\begin{equation}\label{eq:parameters}
 \varepsilon=(1,2,3,4,5),\qquad
 f=\frac1{21}(5,4,4,4,4),\qquad
 g=\frac1{21}(6,6,3,3,3),
\end{equation}
the unique selected set is $S=\{1,2,3,4,5\}$, and its unique
normalized occupation vector is
\begin{equation}\label{eq:occupation}
 \eta=\frac1{100}(14,56,21,8,1).
\end{equation}
There is a nonempty open set of ordered energies and positive
coupling parameters with this same five-mode selection.
Thus $M_S=5>3=2M_B-1$ for $M_B=2$.
\end{theorem}

\begin{proof}
Since $J_3(x)=x^3$, substitution into~\eqref{eq:matrix} gives
$A=B/441$, where
\begin{equation}\label{eq:integer}
 B=\begin{pmatrix}
 0&-16&16&54&128\\
 16&0&-2&-16&-54\\
 -16&2&0&7&56\\
 -54&16&-7&0&7\\
 -128&54&-56&-7&0
 \end{pmatrix}.
\end{equation}
For example, $B_{14}=3^3(5\cdot4-6\cdot3)=54$.
The matrix is skew-symmetric, and direct multiplication yields
\begin{equation}\label{eq:finite}
 B(14,56,21,8,1)^{\mathsf T}=0.
\end{equation}
All entries of~\eqref{eq:occupation} are positive and sum to one.
For the full selected set there are no unselected inequalities
to check in~\eqref{eq:selection}.

To establish stability and rank, let $a$ be any skew-symmetric
$5\times5$ matrix and define its signed Pfaffian vector by
\begin{equation}\label{eq:pfaffian}
 p(a)=\begin{pmatrix}
 a_{23}a_{45}-a_{24}a_{35}+a_{25}a_{34}\\
 -(a_{13}a_{45}-a_{14}a_{35}+a_{15}a_{34})\\
 a_{12}a_{45}-a_{14}a_{25}+a_{15}a_{24}\\
 -(a_{12}a_{35}-a_{13}a_{25}+a_{15}a_{23})\\
 a_{12}a_{34}-a_{13}a_{24}+a_{14}a_{23}
 \end{pmatrix}.
\end{equation}
Then $ap(a)=0$. For its first coordinate, expanding
\begin{align*}
 &-a_{12}(a_{13}a_{45}-a_{14}a_{35}+a_{15}a_{34})\\
 &\quad+a_{13}(a_{12}a_{45}-a_{14}a_{25}+a_{15}a_{24})\\
 &\quad-a_{14}(a_{12}a_{35}-a_{13}a_{25}+a_{15}a_{23})\\
 &\quad+a_{15}(a_{12}a_{34}-a_{13}a_{24}+a_{14}a_{23})
 \end{align*}
gives zero by pairwise cancellation. Relabeling the indices
gives the other four identities. At~\eqref{eq:integer},
\begin{equation}\label{eq:pvalue}
 p(B)=(504,2016,756,288,36)^{\mathsf T}
     =36(14,56,21,8,1)^{\mathsf T}>0.
\end{equation}
The principal minor obtained by deleting the first row and
column has determinant $504^2=254016$, so $B$ has rank at
least four. Every odd-dimensional skew-symmetric real matrix
has zero determinant, since $\det(a)=\det(-a^{\mathsf T})=-\det(a)$.
Consequently $\operatorname{rank}B=4$, and its kernel is one-dimensional.

For the open-set assertion, parametrize the coupling weights by
arbitrary positive $F,G\in\R^5$ through
\begin{equation}\label{eq:normalize}
 f_i=F_i/\textstyle\sum_j F_j,\qquad
 g_i=G_i/\textstyle\sum_j G_j.
\end{equation}
On the open domain of positive $F,G$ and strictly ordered
energies, the entries of $A^{(3)}$ and $p(A^{(3)})$ are continuous.
At the energies in~\eqref{eq:parameters}, take $F=(5,4,4,4,4)$
and $G=(6,6,3,3,3)$.
Then~\eqref{eq:pvalue} gives $p(A^{(3)})=p(B)/441^2>0$.
The intersection of the five inverse images of $(0,\infty)$
under these coordinates is therefore an open neighborhood of
this point. Throughout it,
$\eta=p(A^{(3)})/\sum_i p(A^{(3)})_i$ is a positive kernel vector.
Also the first principal $4\times4$ determinant equals
$p(A^{(3)})_1^2>0$, by expansion of that determinant, and hence
the rank remains four.

Finally, whenever a skew-symmetric $A$ has a strictly positive
kernel vector $\eta$, any solution $\xi$ of~\eqref{eq:selection}
must select every index. Indeed,
$\eta^{\mathsf T}A\xi=-(A\eta)^{\mathsf T}\xi=0$;
if some index were unselected, the left side would be a sum
of nonpositive terms with at least one strictly negative term.
Rank four then makes the normalized full-support vector unique.
This proves all assertions, including uniqueness throughout
the neighborhood.
\end{proof}

\paragraph{Nonnegative transition rates.}
The example fits the algebraic positive-temperature and
population-inverted-bath model of~\cite[Sec.\ V]{ske}.
Put $T_1>0>T_2$, $J_1(x)=x^3$, and $J_2(x)=-x^3$.
For $i\ne j$, set $\Delta=\varepsilon_i-\varepsilon_j$ and
\[
 R_{ij}= f_if_j\frac{\Delta^3}{e^{\Delta/T_1}-1}
       -g_ig_j\frac{\Delta^3}{e^{\Delta/T_2}-1}.
\]
Both summands are positive: in the first fraction the numerator
and denominator have the same sign, and in the second they
have opposite signs before the displayed minus sign.
The identity $(e^x-1)^{-1}-(1-e^{-x})^{-1}=-1$ gives
$R_{ij}-R_{ji}=A^{(3)}_{ij}$; the diagonal may be set to zero.
This is a construction of nonnegative transition rates, not a
claim about an experimental realization of inverted reservoirs.

\begin{corollary}\label{cor:exponents}
There is an open interval of exponents containing $3$ for
which the parameters in~\eqref{eq:parameters} have a unique
five-mode selection. More generally, five-mode selection
holds in a neighborhood of the displayed point in the joint
space $(d,\varepsilon,F,G)$, with $d>0$.
\end{corollary}
\begin{proof}
On the strictly ordered energy domain all off-diagonal energy
differences are nonzero, so $|x|^d\sgn(x)$ is jointly continuous
in $(d,x)$ there. Diagonal entries are identically zero.
Thus the positive signed Pfaffian minors in the preceding
proof remain positive in a joint neighborhood. The same
kernel, rank and uniqueness arguments apply. Intersecting
that neighborhood with the line of fixed energies and
couplings gives the stated interval.
\end{proof}

\begin{corollary}\label{cor:extra}
For every finite $M\geq5$, two cubic product-coupled baths
admit a nonempty open set of positive coupling weights and
ordered energies with exactly five selected modes and a unique
normalized selection vector.
\end{corollary}
\begin{proof}
Let $k=M-5$ and retain the original five energies. Choose distinct
additional energies $E_1,\ldots,E_k>5$. For $t>0$, replace the
first five coupling weights by $f_i/(1+2kt)$ and $g_i/(1+kt)$,
and assign the new weights $2t/(1+2kt)$ and $t/(1+kt)$ to
each extra level. Both weight vectors remain normalized.
The old $5\times5$ principal matrix approaches $A^{(3)}$ as
$t\downarrow0$, so its positive normalized Pfaffian kernel
$\eta(t)$ and rank four persist. Extend $\eta(t)$ by zeros.
For each new row, its drift divided by $t$ tends to
\[
 \sum_{j=1}^5(\varepsilon_j-E_a)^3(2f_j-g_j)\eta_j<0,
\]
because $2f-g=(4,2,5,5,5)/21>0$ and $E_a>5$.
Finitely many strict inequalities therefore hold simultaneously
for all sufficiently small $t>0$, giving exactly the old selected
set by~\eqref{eq:selection}.
If $x,y$ are any two selection vectors, skew symmetry gives
$x^{\mathsf T}Ay=-(Ax)^{\mathsf T}y$; the left side is
nonpositive and the right side nonnegative. Equality forces
their supports to coincide. Rank four then gives normalized
uniqueness. The five positive Pfaffians and the finitely many
negative outside drifts persist in a neighborhood, proving openness.
For $k=0$ this is Theorem~\ref{thm:example}.
\end{proof}

\section{A polynomial extension of the Ohmic bound}
For $M_B\geq1$ baths with a common odd monomial exponent $q$, write
\begin{equation}\label{eq:many}
 A_{ij}=\sum_{b=1}^{M_B}\sigma_b h_i^{(b)}h_j^{(b)}
                 (\varepsilon_j-\varepsilon_i)^q,
 \qquad h_i^{(b)}>0,\quad \sigma_b\ne0.
\end{equation}
The signs of $\sigma_b$ include positive and inverted baths.
Call a selected set nondegenerate if its principal matrix
$A_S$ has a one-dimensional kernel. This holds whenever its
normalized solution of~\eqref{eq:selection} is unique: any
second kernel direction permits a small normalized perturbation
of the strictly positive selected coordinates, preserving the
finitely many strict outside inequalities.

\begin{theorem}\label{thm:bound}
For odd $q\geq1$, every nondegenerate selected set for
\eqref{eq:many} satisfies
\[
 M_S\leq(q+1)M_B-1.
\]
For $q=1$ this recovers the Ohmic bound of Schnell, Ketzmerick
and Eckardt~\cite[Sec.\ V, Eqs.\ (51)--(54)]{ske}.
For two cubic baths it gives $M_S\leq7$.
\end{theorem}
\begin{proof}
Set $n=M_S$ and $K=(q+1)M_B$. On the selected coordinates
form the $n\times K$ matrix $V$ whose columns are
$v_{b,r}=(h_i^{(b)}\varepsilon_i^r)_{i\in S}$,
$0\leq r\leq q$. Let $C$ have entries
\[
 C_{rs}=\begin{cases}(-1)^r\binom qr,&r+s=q,\\0,&r+s\ne q.
 \end{cases}
\]
The binomial expansion gives $A_S=VHV^{\mathsf T}$, where
$H=\operatorname{diag}(\sigma_1C,\ldots,\sigma_{M_B}C)$.
Each $C$ is invertible, since its anti-diagonal entries are
nonzero; it is skew-symmetric because $q$ is odd.
For the positive selected vector $\eta_S$, the vector
$V^{\mathsf T}\eta_S$ is nonzero: its $(b,0)$ entry is
$\sum_{i\in S}h_i^{(b)}\eta_i>0$.
Thus $z=HV^{\mathsf T}\eta_S\ne0$, whereas $Vz=A_S\eta_S=0$.
Consequently $\operatorname{rank}V\leq K-1$ and
$\operatorname{rank}A_S\leq K-1$. A real skew-symmetric
matrix has even rank. Since $K$ is even, its rank is at most
$K-2$. Nondegeneracy gives $n-1=\operatorname{rank}A_S$,
and hence $n\leq K-1$.
\end{proof}

The cubic counterexample and Theorem~\ref{thm:bound} replace
the exponent-independent expectation by the valid
exponent-dependent bound under the stated nondegeneracy
hypothesis. Whether seven can be attained by two cubic baths
is not decided by these arguments. Nondegeneracy matters:
identical oppositely signed couplings give $A=0$, with arbitrary
full-support positive kernel vectors in any dimension.

\begin{corollary}\label{cor:sign}
Suppose the energies are strictly ordered, $d>0$, all couplings
are positive, and all baths have the same spectral sign in
\eqref{eq:many}, with $(\varepsilon_j-\varepsilon_i)^q$
replaced by $J_d(\varepsilon_j-\varepsilon_i)$.
Then precisely one mode is selected: the lowest energy for
positive signs and the highest for negative signs.
\end{corollary}
\begin{proof}
For positive signs, every entry below the diagonal is negative.
If two or more modes were selected, the row of the highest
selected energy would have a strictly negative drift, contrary
to~\eqref{eq:selection}. A singleton other than the lowest
energy gives positive drift in a lower outside row, also a
contradiction. The lowest singleton satisfies all inequalities.
Reversing signs reverses the ordering. This includes the
single-bath equilibrium case in~\cite[Sec.\ II, Eq.\ (12)]{ske}.
\end{proof}

\section{Coverage and verification}
{\small
\begin{longtable}{@{}p{.43\linewidth}p{.50\linewidth}@{}}
\toprule
Source item in arXiv:1802.04088v2 & Mathematical coverage\\
\midrule
Sec.\ V, p.\ 11, Eq.\ (55) and its surrounding paragraph,
non-Ohmic expectation
& Refuted-and-corrected: Theorem~\ref{thm:example} gives an
explicit counterexample; Theorem~\ref{thm:bound} gives the
odd-monomial replacement for nondegenerate selection.\\[4pt]
Sec.\ V, Eqs.\ (51)--(54), Ohmic argument
& Reproved; prior proof Schnell et al., 2018. Theorem~\ref{thm:bound}
recovers $2M_B-1$ at $q=1$ under the explicit nondegeneracy
hypothesis.\\[4pt]
Sec.\ V, Eq.\ (55), odd monomials
& Not posed; exponent-dependent bound proved in
Theorem~\ref{thm:bound}. Its $q=3$ bound is seven, not three,
for two baths.\\[4pt]
Sec.\ V, Eq.\ (55); Schnell (2019), Sec.\ 4.4, p.\ 80,
Eq.\ (4.54), repeated non-Ohmic expectation
& Refuted by the same five-mode example. Corollaries~\ref{cor:exponents}
and~\ref{cor:extra} give exponent neighborhoods and every
finite ambient dimension $M\geq5$.\\[4pt]
Sec.\ II, Eq.\ (11), selection and odd cardinality
& Exactly satisfied; full support is uniquely selected and has
  odd cardinality five.\\[4pt]
Sec.\ II, Eq.\ (12), single-bath equilibrium selection
& Reproved; prior proof Schnell et al., 2018.
Corollary~\ref{cor:sign} also proves the same-sign multi-bath
extension (not posed in this form).\\[4pt]
Knebel et al.\ (2015), Methods, Eq.\ (19)
& Prior proof cited: signed-Pfaffian kernel method.
Equation~\eqref{eq:pfaffian} spells out its five-mode case.\\[4pt]
Sec.\ III A--B, continuum criteria and prefactor near Fig.\ 5;
Sec.\ IV A, random-rate distribution; Sec.\ IV B and Appendix,
kicked-rotor correlations; Sec.\ VI, interacting and photonic systems
& Not covered by this finite product-coupling construction.\\
\bottomrule
\end{longtable}}

The pinned Lean file
\href{https://github.com/the-omega-institute/trureturing/blob/6ffa76e4424ce4c0d1048e56afb0724762ed8285/D5/S3/StatisticalMechanics/BoseSelectionCubicTwoBathRefutation.lean}{\texttt{BoseSelectionCubicTwoBathRefutation.lean}}
contains the declaration \texttt{result} in the namespace
\begin{quote}\small\ttfamily
D5.S3.StatisticalMechanics.\\
BoseSelectionCubicTwoBathRefutation
\end{quote}
Its statement is the negation of \texttt{claim}, which asserts
that the violating set has empty interior for every $d>0$ and
every finite dimension. Its proof exhibits a nonempty open
set at $d=3$, $M=5$, with normalized positive coupling
parameters and at least four selected modes (in fact all five).
The definitions use coordinate-sum normalization
\href{https://github.com/the-omega-institute/trureturing/blob/6ffa76e4424ce4c0d1048e56afb0724762ed8285/D5/S3/Geometry/FiniteCoordinateTruncation.lean}{\texttt{FiniteCoordinateTruncation.normalize}}
and the signed power
\href{https://github.com/the-omega-institute/trureturing/blob/6ffa76e4424ce4c0d1048e56afb0724762ed8285/D5/S3/Analytic/SeriesInequalities/StampachWaclawekExpansionRefutation.lean}{\texttt{StampachWaclawekExpansionRefutation.spow}};
on nonzero real inputs the latter is $x|x|^{d-1}=J_d(x)$,
and its value at zero is zero.
Theorem~\ref{thm:example} above also gives a written proof of this
five-mode construction. Uniqueness, nonnegative rates,
Corollaries~\ref{cor:exponents}, \ref{cor:extra}, \ref{cor:sign}
and Theorem~\ref{thm:bound} have written proofs; these additional
statements are not the scope of the linked declaration.
The accompanying \texttt{verify.py} checks the displayed matrix,
kernel, signed Pfaffian values and principal determinant by
exact rational arithmetic. It checks the five polynomial
kernel identities by coefficient cancellation, rational
six-, seven- and eight-level examples with five selected modes,
and binomial matrix factorizations at $q=1,3,5$.

\begin{thebibliography}{9}\small
\setlength{\itemsep}{3pt}
\bibitem{ske}
A. Schnell, R. Ketzmerick and A. Eckardt,
\emph{On the number of Bose-selected modes in driven-dissipative
ideal Bose gases}, Phys.\ Rev.\ E \textbf{97}, 032136 (2018),
\href{https://doi.org/10.1103/PhysRevE.97.032136}{doi:10.1103/PhysRevE.97.032136};
\href{https://arxiv.org/abs/1802.04088v2}{arXiv:1802.04088v2}.
The posing locator used here is in the arXiv version, Sec.\ V,
p.\ 11, Eq.\ (55) and the surrounding paragraph.
\bibitem{schnell}
A. Schnell, \emph{Dissipative State Engineering in Quantum
Many-Body Systems}, doctoral thesis, Technische Universit\"at
Dresden (2019), Sec.\ 4.4, p.\ 80, Eq.\ (4.54),
\href{https://d-nb.info/1226942288/34}{German National Library full text}.
\bibitem{knebel}
J. Knebel, M. F. Weber, T. Kr\"uger and E. Frey,
\emph{Evolutionary games of condensates in coupled birth-death
processes}, Nat.\ Commun.\ \textbf{6}, 6977 (2015),
\href{https://doi.org/10.1038/ncomms7977}{doi:10.1038/ncomms7977};
\href{https://arxiv.org/abs/1504.07816v1}{arXiv:1504.07816v1},
Methods, Eq.\ (19).
\bibitem{geiger}
P. M. Geiger, J. Knebel and E. Frey,
\emph{Topologically robust zero-sum games and Pfaffian orientation:
How network topology determines the long-time dynamics of the
antisymmetric Lotka--Volterra equation},
Phys.\ Rev.\ E \textbf{98}, 062316 (2018),
\href{https://doi.org/10.1103/PhysRevE.98.062316}{doi:10.1103/PhysRevE.98.062316};
\href{https://arxiv.org/abs/1806.07339}{arXiv:1806.07339}.
\bibitem{petiziol}
F. Petiziol and A. Eckardt,
\emph{Controlling nonequilibrium Bose-Einstein condensation with
engineered environments}, Phys.\ Rev.\ A \textbf{110}, L021701 (2024),
\href{https://doi.org/10.1103/PhysRevA.110.L021701}{doi:10.1103/PhysRevA.110.L021701};
\href{https://arxiv.org/abs/2311.02170}{arXiv:2311.02170}.
\end{thebibliography}
\medskip
\noindent\small\textcopyright\ 2026 Haobo Ma and Wenlin Zhang.
Licensed under \href{https://creativecommons.org/licenses/by/4.0/}{CC BY 4.0}.
\end{document}
